\chapter{Insomnia}
\index{Insomnia}

\section*{Definition and Overview}
Insomnia is difficulty getting to sleep, staying asleep, or waking too early, despite having the opportunity to sleep, combined with daytime effects such as fatigue, poor concentration, and irritability. It is considered chronic when it occurs at least three nights a week for more than three months. Insomnia can be a primary problem or secondary to stress, illness, medicines, or other sleep disorders. The most effective treatment is cognitive behavioural therapy for insomnia (CBT-I), which addresses the thoughts and habits that maintain poor sleep. Sleeping tablets have a limited role and are recommended only for short-term use, because they lose effectiveness and carry risks.

\section*{Epidemiology}
Insomnia is the most common sleep disorder. Around one in three adults has symptoms of insomnia at any time, and around one in ten has chronic insomnia with significant daytime effects. It is more common in women, in older adults, and in people with mental health conditions, chronic pain, and shift work. Insomnia is associated with a substantial burden: more sick days, more accidents, worse work performance, and increased rates of depression, anxiety, and cardiovascular disease. Most people with insomnia do not seek treatment, and of those who do, many are prescribed sleeping tablets rather than receiving the first-line psychological treatment.

\section*{Aetiology and Pathophysiology}
Insomnia develops through a combination of predisposing, precipitating, and perpetuating factors. Predisposing factors include genetics, anxiety proneness, and being a light sleeper. Precipitating factors are the triggers: stress, illness, grief, shift work, or a change in circumstances. Perpetuating factors maintain the problem: worrying about sleep, spending too long in bed, irregular schedules, daytime napping, and using the bedroom for wakeful activities. These behaviours condition the brain to associate bed with wakefulness, so that arousal and sleep do not switch off. The hyperarousal theory of insomnia describes a state of heightened alertness, driven by stress pathways, that prevents normal sleep onset and maintenance.

\section*{Clinical Features}
\begin{itemize}
    \item Difficulty falling asleep, often lying awake for long periods
    \item Waking frequently during the night or too early in the morning
    \item Waking unrefreshed despite adequate time in bed
    \item Daytime fatigue, poor concentration, memory lapses, and irritability
    \item Worrying about sleep, with dread of bedtime
    \item Increased use of caffeine, alcohol, or sleeping tablets
    \item Impaired work, study, and relationships
    \item Physical symptoms such as tension headaches and muscle aches
\end{itemize}

\section*{Differential Diagnosis}
Poor sleep may be due to conditions other than insomnia that need different treatment.
\begin{itemize}
    \item \textbf{Obstructive sleep apnoea:} snoring with pauses in breathing and daytime sleepiness
    \item \textbf{Depression and anxiety:} commonly cause early waking and difficulty getting to sleep
    \item \textbf{Restless legs syndrome and periodic limb movements:} disturbing sleep onset and continuity
    \item \textbf{Circadian rhythm disorders:} delayed or advanced sleep phases, and shift work
    \item \textbf{Medical conditions:} pain, reflux, asthma, thyroid disease, and nocturia
    \item \textbf{Medicines and substances:} caffeine, alcohol, stimulants, and some medicines
    \item \textbf{Other sleep disorders:} narcolepsy and parasomnias
\end{itemize}

\section*{When to Seek Medical Care}
People with insomnia that is persistent, affecting daily function, or accompanied by loud snoring, gasping, or low mood should seek assessment. Review is important before using sleeping tablets regularly, and for insomnia in children or older people that is severe or persistent. A GP can assess causes and refer for CBT-I or a sleep study where needed.

\textbf{Contact your local emergency services for:}
\begin{itemize}
    \item Sudden severe shortness of breath at night, which may indicate heart or lung problems
    \item A person who is very difficult to wake or is breathing abnormally during sleep
    \item Confusion, collapse, or a seizure
    \item Suicidal thoughts or immediate risk of harm
\end{itemize}

\section*{Diagnosis and Investigations}
\begin{itemize}
    \item \textbf{Sleep history:} sleep schedule, patterns, environment, and daytime effects, supported by a sleep diary
    \item \textbf{Screening questionnaires:} the Insomnia Severity Index and Epworth Sleepiness Scale
    \item \textbf{Review of medicines, substances, and lifestyle:} caffeine, alcohol, and medications
    \item \textbf{Mental health assessment:} screening for depression and anxiety
    \item \textbf{Sleep study:} overnight monitoring where sleep apnoea or other disorders are suspected
    \item \textbf{Actigraphy:} a wrist monitor measuring sleep patterns over weeks
    \item \textbf{Blood tests:} thyroid function and other tests as indicated
\end{itemize}

\section*{Management}
\begin{itemize}
    \item \textbf{Cognitive behavioural therapy for insomnia (CBT-I):} the recommended first-line treatment
    \item \textbf{Stimulus control:} using the bed only for sleep and getting up when unable to sleep
    \item \textbf{Sleep restriction:} limiting time in bed to consolidate sleep, then gradually extending it
    \item \textbf{Cognitive therapy:} challenging unhelpful thoughts about sleep
    \item \textbf{Sleep hygiene:} regular schedule, a comfortable environment, and avoiding stimulants
    \item \textbf{Relaxation techniques:} progressive muscle relaxation and mindfulness
    \item \textbf{Short-term medicines:} sleeping tablets for a few weeks only, with caution
    \item \textbf{Treatment of underlying conditions:} apnoea, pain, reflux, and mood disorders
    \item \textbf{Digital programs:} online CBT-I courses, which are effective and accessible
\end{itemize}

\section*{Complications}
\begin{itemize}
    \item Fatigue, poor concentration, and reduced performance at work or study
    \item Increased risk of road and workplace accidents
    \item Depression, anxiety, and irritability
    \item Impaired immunity and increased risk of infection
    \item Associations with obesity, diabetes, high blood pressure, and heart disease
    \item Dependence on sleeping tablets and withdrawal symptoms
    \item Impaired relationships and quality of life
\end{itemize}

\section*{Prognosis and Outcomes}
Chronic insomnia is highly treatable. CBT-I improves sleep in most people and, unlike sleeping tablets, produces lasting change, with benefits that continue after treatment ends. Around two-thirds of people with chronic insomnia respond to CBT-I, and many achieve normal sleep. The earlier treatment is sought, the better the outcome. Because CBT-I is now available online and through GPs and psychologists, effective help is accessible to most people.

\section*{Prevention and Self-Care}
\begin{itemize}
    \item Keep consistent sleep and wake times, including weekends
    \item Get morning light exposure and be active during the day
    \item Avoid caffeine after midday and alcohol near bedtime
    \item Wind down with a relaxing routine before bed
    \item Keep the bedroom cool, dark, and quiet
    \item Get up if you cannot sleep and return when drowsy
    \item Avoid long daytime naps
    \item Put screens away at least 30 to 60 minutes before bed
    \item Seek help early rather than living with poor sleep for years
\end{itemize}

\section*{Sources and Further Reading}
The following reputable sources provide further information for healthcare students and patients seeking reliable, current advice about insomnia.
\begin{itemize}
    \item Sleep Health Foundation. \textit{Insomnia and sleep information.} https://www.sleephealthfoundation.org.au/
    \item Healthdirect Australia. \textit{Insomnia.} https://www.healthdirect.gov.au/
    \item THIS WAY UP (CBT-I online program). \textit{Online insomnia treatment.} https://thiswayup.org.au/
    \item Beyond Blue. \textit{Sleep and mental health.} https://www.beyondblue.org.au/
    \item NHS. \textit{Insomnia.} https://www.nhs.uk/conditions/insomnia/
    \item Mayo Clinic. \textit{Insomnia.} https://www.mayoclinic.org/
\end{itemize}
